% s-016-p016-m004; original PDF 16, 17
\Aphorism{s-016-p016-m004}{7}
Der Mensch ist ein semiotisches Thier; seine Menschheit besteht darin, dass er statt des natürlichen Ausdrucks seiner Bedürfnisse und Befriedigungen sich eine conventionelle, symbolische, nur mittelbar verständliche Zeichensprache angeeignet hat. Er zahlt in Nominalwerthen, in Papieren, das Thier in Baar, in Realwerthen; die Schnelligkeit des Verkehrs, der erleichterte Austausch geistiger Beziehungen hat hierin seinen Grund, und das ist es wohl, was den Menschen zum Menschen gemacht hat? Das Thier thut Ja und Nein, der Mensch sagt Ja und Nein und kommt dadurch wohlfeiler und abstracter zu seinem Glück und um sein Unglück. \Foreign{Ratio} und \Foreign{oratio} sind eine ungeheure Erleichterung des Lebens, und es ist nicht mehr als unparteiisch, wenn im Auftrage der Natur irgend ein Hochmuthsteufelchen den Menschen anspornt, sich durch Schaffung specifisch menschlicher Lust- und Unlustquellen (durch Ideologie) das Leben wieder ein wenig zu erschweren. Oder ist es umgekehrt der erweiterte Horizont und der gesteigerte Widerstand, unter dem der Mensch lebt und leidet, der eine anderweitige Milderung seiner Lebensbedingungen verursacht und rechtfertigt?
\par

